\section*{4. 用 Dwork 方法计算超曲面的 zeta 函数}

\subsection*{4.0.}
设

\begin{equation}\tag{4.0.1}
f(X)=f(X_1,\ldots,X_{n+2})\in\mathbb F_q[X]
\end{equation}

为次数 $d$ 的齐次多项式, 其零点集定义超曲面 $V\subset\mathbb P^{n+1}$.
以 $V^*$ 表示 $V$ 中 $\prod_{i=1}^{n+2}X_i$ 可逆的开子集,
以 $V^*_{\mathrm{aff}}$ 表示 $V^*$ 上的仿射锥.

选取非平凡的加性特征标

\begin{equation}\tag{4.0.2}
\chi:\mathbb F_{q^\nu}\longrightarrow\mu_p(C),
\end{equation}

其中 $C$ 为固定的特征零代数闭域 (例如 $C=\mathbb C$). 因此

\begin{equation}\tag{4.0.3}
\text{对 }x\in\mathbb F_{q^\nu},\qquad
\sum_{y\in\mathbb F_{q^\nu}}\chi(xy)=
\begin{cases}
q^\nu,&x=0,\\
0,&x\ne0.
\end{cases}
\end{equation}

因而对 $x=(x_1,\ldots,x_{n+2})\in(\mathbb F_{q^\nu}^*)^{n+2}$, 有

\begin{equation}\tag{4.0.4}
1+\sum_{x_0\in\mathbb F_{q^\nu}^*}\chi(x_0f(x))
=\sum_{x_0\in\mathbb F_{q^\nu}}\chi(x_0f(x))
=
\begin{cases}
q^\nu,&f(x)=0,\\
0,&\text{其他情形}.
\end{cases}
\end{equation}

在 (4.0.4) 中对 $x\in(\mathbb F_{q^\nu}^*)^{n+2}$ 求和得到

\begin{equation}\tag{4.0.5}
\operatorname{card}V^*_{\mathrm{aff}}(\mathbb F_{q^\nu})
=\frac1{q^\nu}
\sum_{x\in(\mathbb F_{q^\nu}^*)^{n+2}}
\left\{1+\sum_{x_0\in\mathbb F_{q^\nu}^*}\chi(x_0f(x))\right\}.
\end{equation}

由于

\begin{equation}\tag{4.0.6}
\operatorname{card}V^*_{\mathrm{aff}}(\mathbb F_{q^\nu})
=(q^\nu-1)\operatorname{card}V^*(\mathbb F_{q^\nu}),
\end{equation}

公式 (4.0.5) 化为

\begin{equation}\tag{4.0.7}
(q^\nu-1)\operatorname{card}V^*(\mathbb F_{q^\nu})
=\frac{(q^\nu-1)^{n+2}}{q^\nu}
+\frac1{q^\nu}
\sum_{(x_0,x)\in(\mathbb F_{q^\nu}^*)^{n+3}}\chi(x_0f(x)).
\end{equation}

将 $f(X)$ 明确写成单项式之和,

\begin{equation}\tag{4.0.8}
f(X)=\sum_w a_wX^w,
\end{equation}

其中 $w=(w_1,\ldots,w_{n+2})$, $\sum_iw_i=d$,
且 $X^w\mathrel{\mathop:}=X_1^{w_1}\cdots X_{n+2}^{w_{n+2}}$.
于是 (4.0.7) 可改写为

\begin{equation}\tag{4.0.9}
(q^\nu-1)\operatorname{card}V^*(\mathbb F_{q^\nu})
=\frac{(q^\nu-1)^{n+2}}{q^\nu}
+\frac1{q^\nu}
\sum_{(x_0,x)\in(\mathbb F_{q^\nu}^*)^{n+3}}
\prod_w\chi(a_wx_0X^w).
\end{equation}

\subsection*{4.1.}
现在需将“取值于特征零中的加性特征标”具体写出. 自然希望直接取

\begin{equation}\tag{4.1.0}
\chi:x\longmapsto
\exp\!\left(\frac{2\pi i}{p}\,\operatorname{tr}_{\mathbb F_{q^\nu}/\mathbb F_p}(x)\right)
=\xi_p^{\operatorname{tr}_{\mathbb F_{q^\nu}/\mathbb F_p}(x)}.
\end{equation}

\begin{editorialnote}
原书第415页及英文底本的指数中漏去分母 $p$. 上式按 $\xi_p=\exp(2\pi i/p)$ 校正;
迹在指数中取任意整数代表, 所得值与代表的选择无关. 若无此分母, 左端恒为 $1$, 不能给出所需的非平凡加性特征标.
\end{editorialnote}

设

\begin{align}
\tag{4.1.1}
\Omega&=\text{下述域的代数闭包: }\mathbb Q_p,\\
\tag{4.1.2}
\operatorname{ord}:\Omega^*&\longrightarrow\mathbb Q
  &&\text{为赋值, 归一化为 }\operatorname{ord}(p)=1,\\
\tag{4.1.3}
\xi_p&\in\Omega
  &&\text{为本原 $p$ 次单位根.}
\end{align}

可以验证 (参见 [2]), 存在

\begin{equation}\tag{4.1.4}
\pi\in\mathbb Q_p(\xi_p),\qquad
\operatorname{ord}(\pi)=\frac1{p-1},
\end{equation}

使得

\begin{equation}\tag{4.1.5}
\sum_{n\geq0}\frac{\pi^{p^n}}{p^n}=0.
\end{equation}

(也注意到, 左端在 $\mathbb Q_p(\xi_p)$ 中收敛.)

定义 Artin--Hasse 指数

\begin{equation}\tag{4.1.6}
E(X)=\exp\!\left(\sum_{n\geq0}\frac{X^{p^n}}{p^n}\right)
\in\mathbb Q_p[[X]],
\end{equation}

它奇妙地属于 $\mathbb Z_p[[X]]$. 定义

\begin{equation}\tag{4.1.7}
\Theta(X)=E(\pi X).
\end{equation}

其基本性质归纳在下述引理中.

\subsubsection*{引理 4.2.}

\begin{equation}\tag{4.2.1}
\Theta(X)=\sum_n B_n X^n,\qquad B_0=1,\quad B_1=\pi,\quad
\operatorname{ord}(B_n)\geq\frac n{p-1}.
\end{equation}

\begin{equation}\tag{4.2.2}
\Theta(1)=\xi_p,\qquad\text{为本原 $p$ 次单位根.}
\end{equation}

设 $t\in\Omega$, 并假设对某个给定整数 $a>0$, 有

\begin{equation}\tag{4.2.3}
t=t^{p^a};
\end{equation}

即 $t$ 是 $W(\mathbb F_{p^a})\subset\Omega$ 中的 Teichmüller 代表元.
则 $\sum_{j=0}^{a-1}t^{p^j}\in\mathbb Z_p$, 故形式级数

\[
(\Theta(X))^{\sum_{j=0}^{a-1}t^{p^j}}
\]

有定义, 且有形式幂级数恒等式

\begin{equation}\tag{4.2.3.1}
\prod_{j=0}^{a-1}\Theta(t^{p^j}X)
=(\Theta(X))^{\sum_{j=0}^{a-1}t^{p^j}}.
\end{equation}
